\chapter{Conclusion}
\label{ch:conclusion}

\section{Problem and Research Position}

The problem this project addresses is that a single unit of work is recorded in
three tools which version their own records independently, leaving the
relationships between them in people's heads. The decision that fails as a
result is the code review: a reviewer must decide whether one proposed change
still matches the intent it claims to implement, from memory, across several
interfaces, with nothing recording what evidence supported the decision.

The research in Chapter~2 conceded the overlap rather than arguing it away.
Commercial tools already compare a pull request against a linked issue, and one
of them already returns an ``unclear'' outcome instead of forcing a verdict. The
gap that survives that comparison is versioned provenance across tools,
deterministic gating of model output against a stored evidence packet,
reconciliation of asynchronous and imperfectly delivered events, and local
execution. No public corpus labels a design version, a criterion, a diff and a
test together, which is both the reason for this project's own evaluation data
and the reason its claims have to be measured rather than asserted.

\section{Requirements, Design and Implementation to Date}

Chapters~4 and~5 froze the baseline: twenty-nine functional requirements and
fourteen non-functional requirements with an acceptance check attached to each,
a ten-layer architecture with a deployment boundary that only the read-only
adapter and the outbound notification channel cross, a schema in which every
association is a provenance-carrying edge, and a verification design in which
the model is bounded on both sides --- retrieval decides what it may see, and a
deterministic gate decides whether what it produced may be shown.

What has been built in the first iteration is the reviewer surface: nineteen
routes rendered against typed data that encodes the frozen label vocabulary, the
design tokens ported from the exported screens, twelve committed diagrams, the
requirements and design chapters, and the first validated package of evaluation
data. The ingestion, evidence-graph and verification layers described in the
design do not exist yet, and this report does not imply otherwise.

\section{Testing and Validation Status}

The artefacts that exist are checked and they pass: typecheck, lint and a
production build that pre-renders every route; 138 integrity checks over the
fifteen-case dataset pilot with zero failures; one of two blind annotation
sheets complete, with the agreement statistic pending the second. Six defects
are open with named owners, including a test module in the repository that
cannot be collected and the absence of any CI. The claims that carry the
project --- convergence under duplicate delivery, grounded verdicts, abstention
quality, local inference cost --- remain design commitments with acceptance
criteria attached, not measured results.

\section{Risks Carried Forward}

Three risks dominate. The hardware question is unanswered: no model has been
downloaded, so the ceiling on verifier quality is unknown. The reliability and
verification work has not started, which is why the iteration plan in
Chapter~3 is reported with a declared variance instead of a completed status.
And the comparable products have been read but not trialled, so the gap in
Chapter~2 is documented rather than demonstrated under controlled conditions.

\section{Continuation into FYP-2}

The remaining plan in Chapter~10 is sequenced by dependency: contracts and CI
first, then the ingestion reliability layer, then the graph and analyser, then
the packet builder, verifier and gate, with the local-inference spike scheduled
as early as a stable packet exists because its result bounds everything built
after it. The evaluation set is frozen before any component is tuned against it,
the two hard gates are enforced before any metric is reported, and the
adaptation experiment will be reported under a rule written in advance, which
means a negative result is an acceptable outcome.

The project's position at this evaluation is therefore deliberate rather than
defensive: a complete and defensible design, a working reviewer surface, an
evaluated data protocol, and an honest statement of how much of the system
remains to be built.
